% Emerald Publishing - Construction Innovation Submission Template
% by Oleksandr Melnyk
% Ver 0.0.4
% Based on: https://www.emeraldgrouppublishing.com/journal/ci#author-guidelines


\documentclass{article}

\usepackage[english]{babel}

% Set page size and margins
% Replace `letterpaper' with `a4paper' for UK/EU standard size
\usepackage[a4paper,top=2cm,bottom=2cm,left=3cm,right=3cm,marginparwidth=1.75cm]{geometry}

% Useful packages
\usepackage{amssymb}
\usepackage{siunitx}
\PassOptionsToPackage{hyphens}{url}\usepackage{hyperref}
\usepackage{cleveref}
\usepackage[utf8]{inputenc}
\usepackage[right]{lineno}
\usepackage{csquotes}
\usepackage{booktabs}
\usepackage{longtable}
\usepackage{adjustbox}
\usepackage{array}
\usepackage{url}
\usepackage{titlesec}
%\usepackage[compatibility=false]{caption}
\usepackage{authblk}
\usepackage{xcolor} % Load the xcolor package for color options
\renewcommand{\thetable}{\Roman{table}}

% Define a new format for \subsection
\titleformat{\subsection}
  {\mdseries\itshape\large} % Medium series, italic shape, and large font size
  {\thesubsection}{1em}{} % Numbering, spacing, and the section title itself


% Emerald Harvard Citation Style

\usepackage[english]{babel}
\usepackage[style=authoryear,backend=biber,natbib=true,maxcitenames=2,uniquelist=false]{biblatex}
\addbibresource{Bibliography.bib} % your .bib file

% Customizing biblatex for Harvard style
% Customizing biblatex for Harvard style
\DeclareNameAlias{sortname}{family-given}
\DeclareNameAlias{default}{family-given}

\renewbibmacro{in:}{}
\DeclareFieldFormat[article]{title}{\mkbibquote{#1}\addcomma}
\DeclareFieldFormat[book]{title}{\mkbibemph{#1}\addcomma}
\DeclareFieldFormat[bookinbook]{title}{\mkbibemph{#1}\addcomma}
\DeclareFieldFormat[inbook]{title}{\mkbibquote{#1}\addcomma}
\DeclareFieldFormat[incollection]{title}{\mkbibquote{#1}\addcomma}
\DeclareFieldFormat[inproceedings]{title}{\mkbibquote{#1}\addcomma}
\DeclareFieldFormat[manual]{title}{\mkbibemph{#1}\addcomma}
\DeclareFieldFormat[misc]{title}{\mkbibemph{#1}\addcomma}
\DeclareFieldFormat[thesis]{title}{\mkbibemph{#1}\addcomma}
\DeclareFieldFormat[unpublished]{title}{\mkbibquote{#1}\addcomma}
\DeclareFieldFormat[patent]{title}{\mkbibemph{#1}\addcomma}
\DeclareFieldFormat[report]{title}{\mkbibemph{#1}\addcomma}
\DeclareFieldFormat[online]{title}{\mkbibquote{#1}\addcomma}
\DeclareFieldFormat[software]{title}{\mkbibemph{#1}\addcomma}
\DeclareFieldFormat[booklet]{title}{\mkbibemph{#1}\addcomma}
\DeclareFieldFormat[periodical]{title}{\mkbibemph{#1}\addcomma}
\DeclareFieldFormat[standard]{title}{\mkbibemph{#1}\addcomma}

\DeclareFieldFormat[article]{journaltitle}{\iffieldundef{shortjournal}{\mkbibemph{#1}\addcomma}{\mkbibemph{\printfield{shortjournal}}\addcomma}}
\DeclareFieldFormat{volume}{\bibstring{volume}~#1}
\DeclareFieldFormat{number}{\bibstring{number}~#1}

% Definitions for "Vol." and "No."
\DefineBibliographyStrings{english}{
  volume = {Vol.},
  number = {No.}
}

\renewbibmacro*{volume+number+eid}{%
  \printfield{volume}%
  \setunit*{\addspace}%
  \printfield{number}%
  \setunit{\addcomma\space}%
  \printfield{eid}}

\renewbibmacro*{journal+issuetitle}{%
  \usebibmacro{journal}%
  \setunit*{\addcomma\space}%
  \usebibmacro{volume+number+eid}%
  \setunit{\addcomma\space}%
  \usebibmacro{issue+date}}

\renewbibmacro*{publisher+location+date}{%
  \printlist{publisher}%
  \iflistundef{location}
    {\setunit*{\addcomma\space}}
    {\setunit*{\addcolon\space}}%
  \printlist{location}%
  \setunit*{\addcomma\space}%
  \usebibmacro{date}}

\renewcommand*{\bibpagespunct}{\addcomma\space}

% Customizing page field format to prevent duplication
% \DeclareFieldFormat{pages}{%
%   \mkfirstpage[{\mkpageprefix[page]{#1}}]{#1}}

% Customizing citations for Harvard style
\DeclareCiteCommand{\cite}[\mkbibparens]
  {\usebibmacro{prenote}}
  {\usebibmacro{citeindex}%
   \usebibmacro{cite}}
  {\multicitedelim}
  {\usebibmacro{postnote}}

\renewbibmacro*{cite:labelyear+extrayear}{%
  \iffieldundef{labelyear}
    {}
    {\printtext[bibhyperref]{%
       \printfield{labelyear}%
       \printfield{extrayear}}}}

\renewbibmacro*{cite:labeldate+extradate}{%
  \iffieldundef{labelyear}
    {}
    {\printtext[bibhyperref]{%
       \printfield{labelyear}%
       \printfield{extradate}}}}

\AtEveryBibitem{
  \clearfield{month}
  \clearfield{day}
  \ifentrytype{book}{
    \clearlist{location}
  }{}
}

% Formatting "et al." in italics followed by a comma
\DefineBibliographyStrings{english}{
  andothers = {\textit{et al.},}
}

\DeclareFieldFormat[article]{volume}{\bibstring{jourvol}\addnbspace #1}
\DeclareFieldFormat[article]{number}{\bibstring{number}\addnbspace #1}
\DeclareFieldFormat[article]{volume}{Vol. #1}
\DeclareFieldFormat[article]{number}{No. #1}
% Customizing DOI field format to lowercase "doi"
%\DeclareFieldFormat{doi}{\bibstring{doi}\addcolon\space\url{#1}}

% Customizing URL field format to "available at:"
\DeclareFieldFormat{url}{\bibstring{available at}\addcolon\space\url{#1}}
\DeclareFieldFormat{urldate}{\mkbibparens{accessed \addspace#1}}

% Customizing urldate to match the required format
\DeclareFieldFormat{urldate}{%
  \mkbibparens{accessed\space%
    \thefield{urlday}\addspace%
    \mkbibmonth{\thefield{urlmonth}}\addspace%
    \thefield{urlyear}}}


% Configure cleveref
\crefformat{figure}{#2Figure~#1#3}
\Crefformat{figure}{#2Figure~#1#3}
\crefformat{table}{#2Table~#1#3}
\Crefformat{table}{#2Table~#1#3}
\crefformat{section}{#2Section~#1#3}
\Crefformat{section}{#2Section~#1#3}

%Front Matter
\author[1]{Author Name 1}
\author[2]{Author Name 2}

\affil[1]{Institute, Department, City, Country
email (corr. author)}
\affil[2]{Institute, Department, City, Country}

\title{Emerald Publishing Construction Innovation Journal Submission Template}


\begin{document}
\maketitle


\begin{abstract}
\textbf{Purpose} - This study aims to investigate the impact of [specific variable or phenomenon] on [specific outcome or field], addressing a critical gap in the existing literature. The research seeks new insights into [specific problem or issue], contributing to a deeper understanding of [relevant context or application].

\textbf{Design/Methodology/Approach} - The study employs a [qualitative/quantitative/mixed-methods] approach, utilizing [specific methodologies such as surveys, experiments, case studies, etc.]. Data were collected from [description of the sample or data sources] and analyzed using [specific analytical tools or software].

\textbf{Findings} - The results reveal that [briefly describe key findings], indicating that [specific conclusion or implication]. These findings provide evidence that [mention whether the hypothesis was confirmed or if new patterns emerged], offering significant implications for [specific field, industry, or theoretical framework].

\textbf{Originality/Value} - This research offers a novel perspective on [topic], providing valuable insights that extend the current understanding of [subject matter]. The study’s findings contribute to [mention the advancement or practical application], highlighting areas for future research and potential policy or practical applications. 
\\
\\
%add 6 keywords
\textbf{Keywords:} keyword 1; keyword 2; keyword 3, keyword 4; keyword 5; keyword 6
\end{abstract}
\linenumbers


\section{Introduction}
\label{sec:introduction}
Conversational agents (CAs) or chatbots (CBs) are software systems designed to interact with users through natural language, typically using text or voice interfaces. These agents utilize natural language processing (NLP) and machine learning (ML) to comprehend user inputs and provide tailored responses. CAs have become common in customer service, healthcare, and education \citep{araujo2018living}. One key design decision is whether to give these agents human-like characteristics. This practice, known as anthropomorphism, includes features such as human names, conversational language, and personality traits \citep{epley2007seeing}. Anthropomorphic design draws on theories like Computers Are Social Actors (CASA), which suggests that people automatically apply social rules when interacting with computers that display human-like cues \citep{nass1994computers,reeves1996media}. Recent advances in large language models have made it easier to create chatbots that convincingly mimic human conversation \citep{pnas2025anthropomorphic}. A specific form of anthropomorphism involves assigning explicit personas to chatbots, defined sets of attributes that shape their identity and responses. The PersonaChat dataset introduced this approach by providing profile descriptions of conversational agents with characteristics such as hobbies, occupations, and preferences \citep{zhang2018personalizing}. Research on persona-based chatbots has grown since then \citep{chen2024trends}, exploring how these characteristics affect user experience and interaction outcomes.

Despite growing interest in anthropomorphic design for chatbots, we lack a systematic understanding of which persona attributes matter most and their influence user responses. Studies examine individual design characteristics such as names, gender, or communication style \citep{go2019humanizing,seeger2021texting}, but the findings remain scattered across different contexts and applications. Some research suggests that anthropomorphic characteristics increase user trust and satisfaction through mechanisms such as social presence \citep{konya2023anthropomorphism,blut2021understanding}, while other studies identify potential negative effects when human characteristics create unrealistic expectations \citep{crolic2022blame,ciechanowski2019shades}. The field needs synthesis of empirical evidence to guide design decisions for information systems that use conversational agents. This meta-analysis addresses this gap by systematically reviewing studies on persona-based chatbots to identify which anthropomorphic characteristics have been studied, how they are implemented and what effects they produce on user outcomes.

Chatbot personas can encompass a wide range of human-like attributes. Some chatbots have demographic characteristics, such as age and gender, while others possess personality traits or communication styles \citep{go2019humanizing}. Visual features like avatars or profile pictures can make chatbots appear more human \citep{araujo2018living}. Conversational design choices include whether the chatbot uses a human name, refers to itself in first person, or includes backstory information \citep{seeger2021texting}. Understanding which attributes researchers have examined helps to identify what design elements matter for creating effective conversational agents \citep{blut2021understanding}. A systematic inventory of persona attributes also reveals gaps where certain characteristics remain understudied. Recent work has explored how different persona types affect user perceptions \citep{sun2024building} and how personas can be implemented in modern large language model systems \citep{chen2024trends}. \textbf{RQ1:} \textbf{What types of persona attributes have been examined in research on chatbot design?}

The implementation choices shape the way users experience chatbot personas. Researchers must decide whether to make personas explicit or implicit, fixed or adaptive to individual users \citep{huang2023personalized}. The decision to disclose the artificial nature of a chatbot versus present it as human-like raises ethical considerations \citep{araujo2018living}. The technical implementation varies between platforms, from text-only interfaces to embodied avatars with voice capabilities \citep{go2019humanizing}. Some studies provide users with written persona descriptions before interaction \citep{zhang2018personalizing,dinan2020second}, while others embed persona characteristics in the conversation itself \citep{song2021bob}. Recent work explores how personas can be maintained across multiple conversation sessions \citep{xu2022beyond} and how they can be grounded to specific user queries \citep{jang2022focus}. These methodological variations affect what conclusions we can draw about the effectiveness of personas. \textbf{RQ2:} \textbf{How are persona attributes operationalized and implemented in chatbot systems?}

The primary motivation for studying chatbot personas is understanding their effects on user experience and behavior. Research examines various outcome measures including trust, satisfaction, perceived intelligence, and behavioral intentions \citep{konya2023anthropomorphism,blut2021understanding}. Some studies focus on immediate interaction quality \citep{go2019humanizing}, while others investigate the building of long-term relationships \citep{xu2022beyond}. The effects may differ according to the application context, and educational chatbots require different design choices than customer service agents \citep{araujo2018living}. User characteristics like emotional state can moderate how people respond to anthropomorphic features, with angry customers showing negative reactions to human-like chatbots \citep{crolic2022blame}. Individual differences in personality and cognitive style also influence the perceptions of anthropomorphism \citep{spatola2021personality}. Synthesizing evidence on which persona attributes produce which effects under what conditions provides actionable guidance for designers. \textbf{RQ3:} \textbf{What effects do different persona attributes have on user perceptions, attitudes, and behaviors?}

To address these research questions, we conducted a systematic review and meta-analysis of empirical studies examining persona-based chatbots. We searched multiple academic databases to identify studies published between 2018 and 2024 that investigated chatbots with explicit persona attributes. Our search focused on studies that operationalized personas through specific characteristics like demographic features, personality traits, or communication styles. We systematically coded each study for the types of persona attributes examined, implementation methods used, study contexts, and outcome variables measured. For studies reporting quantitative effect sizes, we extracted statistical information to enable meta-analytic synthesis. Our coding scheme captured persona attributes across multiple dimensions including visual representation, conversational features, demographic characteristics, and personality traits. We also recorded methodological details such as whether personas were disclosed to users, whether they were fixed or adaptive, and the application domain. This systematic approach allows us to identify patterns across studies and provide comprehensive answers to our research questions about what persona attributes exist in the literature, how they are implemented, and what effects they produce.

\section{Method}
\label{sec:method}
The \textcolor{teal}{\textit{Methods}} section should provide a comprehensive and detailed description of the research procedures employed in the study. 
It should explain the data collection and analysis techniques, ensuring that the description is thorough enough to allow others to replicate the study. 
The section should be written in the past tense and typically employ passive voice to maintain an objective tone. 
Additionally, it is essential to include all necessary bibliographic information for each source cited, such as in the example \cite{fellows_research_2021}.

Use the Emerald Harvard style defined in this template for references, and adjust the citation style if necessary. Include DOIs where available. 

\section{Main Section}
\label{sec:main_section}
When submitting manuscripts using \LaTeX, a PDF file must also be included. 
Articles should be between 6,000 and 10,000 words, encompassing all text, tables, and figures. 
A concise and descriptive title must be included, and it is essential to list all contributing authors in the submission, along with their email addresses, names, and affiliations.



\subsection{Subsection 1}
Ensure headings are concise and clearly indicate the hierarchy. 
Use sparingly and identify with consecutive numbers in square brackets.
Submit figures (such as \cref{fig-1} electronically at the highest resolution. Number figures consecutively with clear captions.

\begin{figure}[ht]
 \centering
 \makebox[\textwidth][c]{\includegraphics[width=1\textwidth]{Figures/Figure_1.jpg}}%
 \caption{Figure caption}
 \label{fig-1}
\end{figure}


\subsection{Subsection 2}

\textbf{Emerald Publishing} accepts formats such as .ai, .eps, .jpeg, .bmp, and .tif. Electronic figures created in other applications should be provided in their original formats. 
Additionally, these figures should either be copied and pasted into a blank MS Word document or submitted as a PDF file.

\section{Results}
\label{sec:modelling}
This \LaTeX template is designed to create an academic article with specific formatting and citation requirements. Here's a plain-text explanation of how the template has been structured:

Useful Packages:
        Several \LaTeX packages are included to provide additional functionalities:
            \begin{itemize}
                \item amssymb: Provides access to additional mathematical symbols.
                \item siunitx: Allows for consistent formatting of numbers and units.
                \item hyperref: Adds hyperlinks to the document, with an option to handle URL line breaks.
                \item cleveref: Enhances cross-referencing features.
                \item inputenc: Handles input encoding, using UTF-8 for wider character support.
                \item lineno: Adds line numbers to the document.
                \item csquotes: Handles quotation marks in a language-sensitive way.
                \item booktabs: Provides better formatting for tables.
                \item longtable: Supports tables that can span multiple pages.
                \item adjustbox: Allows for resizing and adjusting the positioning of content.
                \item array: Provides additional tools for formatting tables and arrays.
                \item url: Simplifies the inclusion of URLs in the document.
                \item titlesec: Customizes section titles (though the specific customization is commented out).
                \item authblk: Manages author affiliations.
            \end{itemize}


Submit tables (such as \ref{tab-II}) separately with clear titles and label their positions in the main text.\\
Non-critical supplementary material should be submitted separately and will be hosted alongside the article.

\begin{table}[ht]
 \centering
 \makebox[\textwidth][c]{\includegraphics[width=1\textwidth]{Figures/Table_I.jpg}}%
 \caption{Tables numbered in roman numerals}
 \label{tab-II}
\end{table}

The table numbering format is redefined with the command below. This command switches from Arabic to Roman numerals for table captions:
        \begin{verbatim} 
        \renewcommand{\thetable}{\Roman{table}} 
        \end{verbatim}


\section{Discussion}
\label{sec:discussion}
A research paper's \textcolor{teal}{\textit{Discussion}} section typically encompasses several elements. 
First, it involves the interpretation of results, where you explain what your findings signify in relation to your research question. 
This is followed by a comparison with previous work, which involves relating your results to existing literature and highlighting similarities and differences. 
If applicable, you should also explain any unexpected results, discussing surprising findings and their potential causes. 
Additionally, this section should address the broader implications of your results for the field or related areas. 
Acknowledging any limitations of your study or methodology is crucial, as is suggesting potential avenues for future research based on your findings.

\section{Conclusion}
\label{sec:conclusion}

A strong \textcolor{teal}{\textit{Conclusion}} typically encompasses several essential elements. It begins with a restatement of the thesis or main idea, reinforcing the core message of the paper. This is followed by a summary of the key points discussed throughout the text, which helps to consolidate the main arguments. The conclusion should also include final thoughts or reflections on the topic, providing a thoughtful discussion wrap-up. 

\hfill

\textit{Acknowledgments} \\Mention all external funding sources in the acknowledgements.\\

\textit{Disclosure statement:} \\No potential conflict of interest is reported by the authors




%\break
\printbibliography

\end{document}
